\chapter{Evaluation}
\label{chap:evaluation}

This chapter defines the cross-domain use case used to demonstrate AgenTwin,
    followed by the evaluation methodology that assesses the artifact's
    semantic, functional, and performance characteristics.
The use case provides the concrete operational scenario against which all evaluation
    dimensions are measured, grounding the assessment in a realistic deployment context.

\section{Cross-Domain Use Case Definition}
\label{sec:use-case}

To validate AgenTwin's agent-mediated interoperability approach,
    a realistic cross-domain scenario has been designed that requires
    cooperation between Smart City and Energy Grid DT systems.

The use case addresses the central operational challenge of public transport depot charging:
    assigning overnight charging stations to a mixed electric bus fleet
    while respecting the depot's site-wide power limit imposed by
    the distribution transformers serving the depot area.
In public transport depot charging,
    the overnight cycle is the operational backbone of electric bus service:
    it determines daily fleet readiness and directly affects total operating cost.
Vehicle duty cycles and the required daily energy are modelled per route and per bus type,
    meaning that resolving vehicle models to their electrical specifications
    is a prerequisite for any grid-safe assignment.
Critically, the depot's maximum import capacity sets a site-wide power limit
    that constrains how much power can be drawn simultaneously across all charging bays,
    requiring that assignments respect transformer zone headroom
    rather than treating all stations as interchangeable~\cite{elintacharge2026}.

\textbf{User Persona:} % based on Hallcon job post for Coordinator, EV Charging Operations: https://talents.vaia.com/companies/hallcon/coordinator-ev-charging-operations-22312420
An EV Operations Coordinator manages the day-to-day readiness of a transit vehicle fleet,
    overseeing vehicle and driver activity, communicating operational changes,
    and reporting deficiencies across the team.
Each night, as buses return to depot, the coordinator must assign charging stations
    to a mixed fleet whose individual power requirements vary by vehicle model.
The coordinator opens AgenTwin and submits a natural-language request such as:
    \textit{``I need to charge 25 Mercedes-Benz eCitaro and 25 Caetano e.City Gold buses overnight.''}

\textbf{Cross-Domain Requirement:}
This request exposes a known industry challenge in depot charging:
    interoperability issues across vehicle OEMs,
    charging interfaces, and backend systems~\cite{elintacharge2026}.
In AgenTwin's architecture, this manifests as a split between two isolated platforms:

\begin{itemize}
    \item The Smart City DT knows charging station locations, slot availability,
        operational status, and supported vehicle types,
        but has no knowledge of transformer loads or grid capacity constraints
    \item The Energy Grid DT monitors transformer loads and computes available headroom per zone,
        but does not model charging station infrastructure, slot availability, or vehicle types
\end{itemize}

AgenTwin bridges these two platforms:
    it resolves vehicle models to their charge rate requirements,
    correlates charging station locations with grid topology zones,
    and produces a complete assignment that satisfies both service availability
    and the depot's maximum import capacity constraint.

\textbf{Interaction Workflow:}
\begin{enumerate}
    \item \textbf{Operator submits request} via natural language through user interface,
        identifying vehicle models and fleet size.
    \item \textbf{Vehicle enrichment:} the Energy Grid agent resolves each model name against
        \texttt{ElectricVehicle} entities in Ditto,
        retrieving \texttt{chargeRateMax} per model without requiring
        the operator to provide electrical specifications.
    \item \textbf{Infrastructure query:} the Smart City agent queries \texttt{EVChargingStation}
        entities in KTWIN, filtering by
        \texttt{availableCapacity > 0},
        \texttt{status = working}, and
        \texttt{allowedVehicleType = bus}.
    \item \textbf{Grid capacity query:} the Energy Grid agent aggregates
        \texttt{CustomerLoad.p} values per transformer zone in Ditto,
        computing available headroom as
        $\texttt{headroom} = \texttt{ratedS} - \Sigma\texttt{CustomerLoad.p}$.
    \item \textbf{Geographic--topological bridge:} the agents correlate each candidate station's
        GeoJSON coordinates (\texttt{GeoLocation}) with the CIM \texttt{Location} position points
        of transformer zones to determine which transformer serves each station.
        This step constitutes the primary semantic gap: the two platforms
        represent geographic position in structurally incompatible formats with no shared identifier.
    \item \textbf{Assignment computation:} each bus is mapped to a station whose transformer zone
        can absorb its \texttt{chargeRateMax} while maintaining at least 5\% safety headroom.
        Bus models with higher power requirements are assigned to zones with greater available capacity.
    \item \textbf{Assignment proposal presented} to the operator as a structured natural-language
        response listing each bus and its assigned charging station.
\end{enumerate}

\textbf{Evaluation Scenario:}
Both DT platforms are pre-populated with a known depot state:
    transformer zones at representative overnight baseline loads and
    charging stations reflecting realistic slot availability.
The evaluation submits the bus fleet request and verifies that
    AgenTwin produces a complete, constraint-satisfying assignment
    without any manually coded mapping between the two platforms.

\textbf{Success Criteria:}
\begin{itemize}
    \item All 50 buses receive a valid charging station assignment
    \item No transformer zone exceeds 95\% of rated capacity after the full assignment
    \item Only bus-capable stations (\texttt{allowedVehicleType = bus}) are assigned
\end{itemize}

\subsection{Semantic Gap Characterization}
\label{subsec:semantic-gaps}

The scenario requires bridging multiple semantic gaps between the
    \texttt{smart-city} ontology (KTWIN) and the
    \texttt{energy-grid} ontology (Ditto).
These gaps are representative of real-world DT integration challenges:
    different domains model the same physical infrastructure from different perspectives using
    domain-specific ontologies optimized for their use cases.
Table~\ref{tab:semantic-gaps} summarizes the main semantic gaps identified.

\begin{table}[H]
    \centering
    \small
    \begin{tabular}{ | m{4cm} | m{4cm} | m{5cm} | }
        \hline

        \textbf{Smart City Concept} & \textbf{Energy Grid Concept} & \textbf{Semantic Challenge} \\
        \hline

            \texttt{EVChargingStation}

            (physical infrastructure asset)
            
            &
            
            \texttt{CustomerLoad}

            (grid load entity)
            
            &
            
            Same physical installation, no shared identifier; agent must correlate by geographic proximity. \\
        \hline
            
            \texttt{GeoLocation}
            
            (GeoJSON,
            WGS84 implicit per RFC~7946)
            
            &
            
            \texttt{Location},
            
            \texttt{PositionPoints},
            
            \texttt{CoordinateSystem}

            (CRS declared as OGC URN)
            
            &

            Different coordinate reference systems with different structural encodings.\\
        \hline

    \end{tabular}
    \caption{Semantic gaps requiring agent-mediated mapping}
    \label{tab:semantic-gaps}
\end{table}

\section{Evaluation Methodology}
\label{sec:evaluation-methodology}

The evaluation assesses whether AgenTwin achieves its design goal:
    enabling agent-mediated semantic interoperability between heterogeneous Digital Twin systems
    without requiring manually coded integration between platforms.
Following the DSR evaluation principle of demonstrating artifact feasibility~\cite{peffers2007},
    the scope is deliberately bounded to the fleet charging coordination use case.

The evaluation is structured as three layered dimensions that build on one another.
Semantic mapping evaluation establishes, at the component level, that agents correctly
    bridge each identified semantic gap before any end-to-end claim is made.
Functional evaluation then verifies, end-to-end, that the artifact produces correct
    outputs across both nominal and boundary operating conditions.
Performance evaluation characterises the operational cost of the integration mechanism,
    informing whether the approach is practical for deployment.

\subsection{Semantic Mapping Evaluation}
\label{subsec:semantic-mapping-evaluation}

Semantic mapping evaluation assesses whether agents correctly bridge each semantic gap
    identified in Table~\ref{tab:semantic-gaps}.
To ensure test/production parity, the agent is designed to always emit a
    structured mapping declaration as part of its response.
This artifact names the entities correlated, the reasoning used,
    and any coordinate or encoding transformations applied.

\begin{itemize}
    \item \textbf{Mapping Coverage:} % Suggestion: use deepEval JSONCorrectnessMetric
    Verify that the structured mapping declaration contains a valid entry for each
        semantic gap in Table~\ref{tab:semantic-gaps}.
    Evaluated via a JSON correctness metric that verifies whether the structured output
        conforms to the expected schema: expected gaps must exist or the test fails.

    \item \textbf{Mapping Correctness:} % Suggestion: use deepEval DAGMetric (BinaryJudgementNode per gap)
    Verify that each entry in the mapping declaration is semantically accurate
        against the specific challenge of its gap.
    Evaluated via a deterministic decision tree that uses LLM-as-a-judge at each
        branching point:
    \begin{itemize}
        \item \texttt{EVChargingStation} $\leftrightarrow$ \texttt{CustomerLoad}:
            the agent must have correlated entities by geographic proximity.
        \item \texttt{GeoLocation} $\leftrightarrow$
            \texttt{Location}:
            the agent must have correctly identified and resolved the structural
            difference between WGS84 GeoJSON and OGC URN-declared
            coordinate reference system.
    \end{itemize}
\end{itemize}

\subsection{Functional Evaluation}
\label{subsec:functional-evaluation}

Functional evaluation verifies that the artifact produces correct outputs
    across two operating points: a nominal scenario where transformer capacity is sufficient
    to accommodate the full fleet, and a stress scenario where it is not.
Together, these two cases confirm that the system satisfies hard constraints
    regardless of whether demand fits within available headroom.

\begin{itemize}
    \item \textbf{Nominal Scenario:} % Suggestion: use deepEval TaskComplitionMetric + programmatic assertions
    Execute the cross-domain use case described in Section~\ref{sec:use-case}
        with transformer headroom sufficient for all 50 buses.

    \item \textbf{Oversubscription Scenario:} % Suggestion: use deepEval TaskComplitionMetric (strict_mode=True)
    Reduce available transformer headroom so that it cannot accommodate the full fleet.
    Verify that the system aborts the assignment and explains why it could not proceed
        (e.g., which transformer zone would be exceeded).
\end{itemize}

\subsection{Performance Evaluation}
\label{subsec:performance-evaluation}

Performance evaluation therefore characterizes AgenTwin's operational overhead
    against the constraints of the deployment context itself,
    and presents the results as an open discussion of practical feasibility.
The goal is to determine whether the approach imposes prohibitive costs
    for realistic operational use,
    with a manually-coded integration baseline identified as a direction for future work.

\begin{itemize}
    \item \textbf{Task Completion Time:} % Suggestion: use LangFuse
    Measure total time from operator submission to delivery of the complete assignment proposal,
        including any clarification turns between the agent and the operator.
    Report mean, median, and 95th percentile times across multiple executions.
    Results are discussed against the operational constraint that the assignment must be
        ready before the charging cycle begins to ensure fleet readiness by scheduled
        departure time.

    \item \textbf{LLM Resource Consumption:} % Suggestion: use LangFuse
    Record the number of LLM inference calls and total tokens consumed
        per assignment request.
    These metrics constitute the dominant cost driver in agent-mediated systems
        and determine whether the approach is deployable under realistic budget constraints.
    Results are reported descriptively to inform infrastructure requirements
        and support future cost comparisons as the research matures.
\end{itemize}
